% !TeX root = ../navier-stokes-manuscript.tex
\subsection{The Sobolev column inside every temporal row}\label{sub:SobolevRegularityInsideEveryCompletedRow}

The critical-height table was built from the base velocity \(V_0=u\).
The mixed Tower asks for more: every fixed response \(V_r\) has to carry its own spatial column.
The equation gives exactly that.

Begin with the first nonzero row.
Equation \eqref{eq:BodySecondTemporalVelocityRow} can be written as
\[
  V_2
  +(V_1\cdot\nabla)V_0
  +(V_0\cdot\nabla)V_1
  +\nabla Q_1
  =\nu\Delta V_1.
\]
For \(s\geq0\), define
\[
  E_{1,s}:=\frac12\norm{\Lambda^sV_1}_2^2
\]
and
\[
  \mathcal P_{1,s}
  :=-
  \operatorname{Re}
  \pair{\Lambda^sV_1}
  {\Lambda^s\!\left((V_1\cdot\nabla)V_0+(V_0\cdot\nabla)V_1+\nabla Q_1\right)}_{L^2}.
\]
Pairing this row with \(\Lambda^{2s}V_1\) gives
\begin{equation}\label{eq:BodyFirstTemporalRowSpatialEnergy}
  E_{1,s}'=\mathcal P_{1,s}-2\nu E_{1,s+1}.
\end{equation}
The first temporal response has the same viscous spatial ascent as the base velocity.

The calculation is unchanged at an arbitrary fixed temporal order.
Set
\begin{align*}
  E_{r,s}
  &:=\frac12\norm{\Lambda^sV_r}_2^2,
  \\
  \mathcal P_{r,s}
  &:=-\operatorname{Re}
  \pair{\Lambda^sV_r}
  {\Lambda^s\!\left(
    \sum_{a=0}^{r}\binom{r}{a}(V_a\cdot\nabla)V_{r-a}
    +\nabla Q_r
  \right)}_{L^2}.
\end{align*}
Using \eqref{eq:BodyTemporalVelocityRecurrence},
\begin{equation}\label{eq:BodyEveryTemporalRowSpatialEnergy}
  E_{r,s}'=\mathcal P_{r,s}-2\nu E_{r,s+1}.
\end{equation}
Starting at \(s=1/2\) and repeating the completed-height calculation yields
\begin{equation}\label{eq:BodyEveryTemporalRowCompletedHeight}
  \partial_t^nE_{r,1/2}
  =(-2\nu)^nE_{r,n+1/2}
  +\sum_{j=0}^{n-1}
  (-2\nu)^j
  \partial_t^{\,n-1-j}\mathcal P_{r,j+1/2}.
\end{equation}
Viscosity does not identify temporal order with spatial order.
It opens the same spatial ladder separately inside each selected temporal response.

Equation~\eqref{eq:BodyEveryTemporalRowCompletedHeight} identifies the spatial ladder; it does not yet give a pointwise-in-time bound.
That bound appears when a whole-terminal rectangle supplies the adjacent pair
\[
  V_r\in L^2(I_*;H^{m+1}),
  \qquad
  V_{r+1}=\partial_tV_r\in L^2(I_*;H^{m-1}).
\]
The trace identity then gives
\[
  \sup_{0\leq t\leq T_*}\norm{V_r(t)}_{H^m}<\infty.
\]
In three dimensions,
\[
  H^m(\R^3)\hookrightarrow C_b^k(\R^3)
  \qquad\text{when}\qquad
  m>k+\frac32.
\]
For any fixed spatial derivative order \(k\), choosing a finite integer \(m>k+3/2\) turns that adjacent trace into pointwise control of every derivative of \(V_r\) through order \(k\).

The two independent directions can now be read in one grid.

\begin{center}
\captionsetup{hypcap=false}
\captionof{table}{The mixed derivative grid. A finite cell records one membership \(V_r\in H^m\); temporal order runs downward and spatial height runs across.}
\label{tab:MixedDerivativeGrid}
\renewcommand{\arraystretch}{1.4}
\begin{tabular}{@{}c c c c c c@{}}
\toprule
& \(H^0_x\) & \(H^1_x\) & \(H^2_x\) & \(H^3_x\) & \(\cdots\) \\
\midrule
\(V_0=u\) & \(V_0\in H^0\) & \(V_0\in H^1\) & \(V_0\in H^2\) & \(V_0\in H^3\) & \(\cdots\) \\
\(V_1=u_t\) & \(V_1\in H^0\) & \(V_1\in H^1\) & \(V_1\in H^2\) & \(V_1\in H^3\) & \(\cdots\) \\
\(V_2=u_{tt}\) & \(V_2\in H^0\) & \(V_2\in H^1\) & \(V_2\in H^2\) & \(V_2\in H^3\) & \(\cdots\) \\
\(V_3=u_{ttt}\) & \(V_3\in H^0\) & \(V_3\in H^1\) & \(V_3\in H^2\) & \(V_3\in H^3\) & \(\cdots\) \\
\(\vdots\) & \(\vdots\) & \(\vdots\) & \(\vdots\) & \(\vdots\) & \(\ddots\) \\
\bottomrule
\end{tabular}
\end{center}

Now read one row in both directions.
A finite singular time forces
\[
  \sup_{0<t<T_*}\norm{V_r(t)}_3=\infty
\]
by \eqref{eq:TemporalTowerMustEscape}.
If a whole-terminal adjacent pair centred at \(H^1\) is available, then \(V_r\in C([0,T_*];H^1)\), and hence
\[
  \sup_{0<t<T_*}\norm{V_r(t)}_3
  \leq
  C\sup_{0<t<T_*}\norm{V_r(t)}_{H^1}
  <\infty.
\]
This is the collision the whole-terminal theorem will create inside the same temporal row.
The response that a finite-time singularity asks to lose control is also a row in which viscosity opens stronger spatial regularity.

The identities alone do not provide the needed bounds.
They identify the object the estimate has to control: every selected finite temporal depth together with every selected finite spatial height.
At the terminal time, the smallest useful piece is not one cell but an adjacent pair that produces a strong middle-space trace.
